\Needspace{18\baselineskip}
\section{Worked problems}

\Needspace{12\baselineskip}
\subsection{Prediction: address translation}
\textbf{Problem.} \emph{Prediction.} With $B=4$ and $T_A=[7,2,11]$, identify the physical
block and in-block offset for logical positions 0, 3, 4, 5 and 10.

\textbf{Solution.} With $B=4$, logical positions 0 and 3 are in logical block 0, so they map to
physical block 7 at offsets 0 and 3. Positions 4 and 5 are in logical block 1,
so they map to physical block 2 at offsets 0 and 1. Position 10 is in logical
block 2 at offset 2, so it maps to physical block 11, offset 2.

\Needspace{12\baselineskip}
\subsection{Derivation: exact fragmentation}
\textbf{Problem.} \emph{Derivation.} For lengths $(1,4,5,8,9)$ and block size four,
compute each $n_i$, each tail waste $w_i$, total paged utilization and the
worst-case bound on total tail waste.

\textbf{Solution.} For lengths $(1,4,5,8,9)$ and $B=4$, block counts are $(1,1,2,2,3)$ and tail
wastes are $(3,0,3,0,3)$. Valid positions total 27; allocated positions total
$9\times4=36$, so utilization is $27/36=0.75$. With five live nonempty
requests the strict bound is total tail waste $<5B=20$ positions; the tighter
integer bound is at most $5(B-1)=15$. Actual waste here is 9.

\Needspace{12\baselineskip}
\subsection{Bytes: from one KV scalar to a hundred blocks}
\textbf{Problem.} \emph{Bytes.} For $L=32$, $H_{kv}=8$, $d_h=128$, $b_{kv}=2$ bytes and
$B=16$, derive bytes per token, bytes per physical block across all layers,
and payload bytes represented by a table with 100 blocks.

\textbf{Solution.} One layer stores
$2H_{kv}d_hb_{kv}=2(8)(128)(2)=4096$ bytes per token. Across 32 layers that is
131,072 bytes = 128 KiB per token. A 16-token block therefore carries
2,097,152 bytes = 2 MiB of KV payload across the model. One hundred such
blocks represent 200 MiB of payload (209,715,200 bytes), excluding allocator
metadata and alignment.

\Needspace{12\baselineskip}
\subsection{Correctness: shared partial tail}
\textbf{Problem.} \emph{Correctness.} Two requests share a partial tail with reference
count two. Explain why incrementing the reference count alone does not make an
append safe, and state the required COW transition.

\textbf{Solution.} A reference count proves that more than one owner can read the block; it does
not isolate writes. Before X appends, X allocates a private block, copies the
valid cells of the shared tail, changes its final table entry to the new
block, decrements the old block's reference count, and writes only after the
new block is private. Y continues to observe the old bytes.

\Needspace{12\baselineskip}
\subsection{Falsification: paging can lose on a chosen workload}
\textbf{Problem.} \emph{Falsification.} Construct a workload for which a chosen paged
block size consumes more physical token capacity than a tightly sized
contiguous allocation. State exactly which benefit of paging your example
does and does not falsify.

\textbf{Solution.} Let every request have exactly five tokens and compare a tightly sized
contiguous allocation of five positions per request with paging at $B=8$.
Paging allocates eight positions per request and therefore consumes 60\% more
token capacity. This falsifies ``paging always uses less memory than every
contiguous layout.'' It does not falsify paging's advantage over
maximum-length reservation or its ability to avoid external fragmentation and
support non-contiguous growth.
